% Shared style of the LearningToOptimize diagrams: the palette, lane frames, cards,
% chips and glyphs of the paper-style TikZ figures (the same visual system as the
% PortfolioOpt.jl and DecisionRules.jl figures). Each figure \input's this file.
\usepackage{amsmath,amssymb,bm}
\usepackage[T1]{fontenc}
\usepackage[scaled=0.92]{helvet}
\renewcommand{\familydefault}{\sfdefault}
\usetikzlibrary{arrows.meta,shapes.geometric,shapes.symbols,calc,positioning,decorations.pathreplacing}

\definecolor{dTeal}{HTML}{1A9AA0}
\definecolor{dTealLine}{HTML}{3FB3B8}
\definecolor{dNavy}{HTML}{3F4C8C}
\definecolor{dInk}{HTML}{1E2235}
\definecolor{dRed}{HTML}{E4573D}
\definecolor{dCoral}{HTML}{E2765A}
\definecolor{dOrange}{HTML}{F4A93B}
\definecolor{dPeach}{HTML}{FBD9AC}
\definecolor{dPanel}{HTML}{EEF0F2}
\definecolor{dGray}{HTML}{8A90A0}
\definecolor{dNode}{HTML}{14A3A8}
\definecolor{dShadow}{HTML}{D5DAE3}

% \ifcompact is true when the job name ends in "-compact" (narrow layout).
\newif\ifcompact
\ifnum\pdfmatch{-compact$}{\jobname}=1 \compacttrue\fi

\tikzset{
  >={Stealth[length=5pt,width=4pt]},
  every node/.style={text=dInk},
  title/.style={anchor=west,font=\small,inner sep=0pt},
  head/.style={font=\small\bfseries,text=dNavy},
  note/.style={font=\itshape\footnotesize,text=dNavy},
  dnote/.style={font=\itshape\footnotesize,text=dCoral},
  code/.style={font=\ttfamily\footnotesize,text=dNavy},
  store/.style={cylinder,shape border rotate=90,aspect=0.28,fill=dPeach,draw=dOrange,
                line width=0.7pt,minimum width=0.95cm,minimum height=1.2cm,font=\normalsize},
  dia/.style={diamond,fill=dOrange,draw=none,minimum size=0.9cm,inner sep=0pt,aspect=1,
              font=\normalsize},
  lnode/.style={circle,fill=dRed,draw=none,minimum size=0.66cm,inner sep=1pt,
                text=white,font=\normalsize},
  card/.style={rectangle,rounded corners=4pt,fill=white,draw=dShadow,line width=0.6pt},
  sheet/.style={rectangle,rounded corners=3pt,fill=dPanel,draw=dShadow,line width=0.6pt},
  chip/.style={rectangle,rounded corners=3pt,fill=dPanel,draw=none,font=\small,
               inner xsep=4pt,inner ysep=3pt},
  ochip/.style={chip,fill=white,draw=dGray,dashed,line width=0.6pt},
  pchip/.style={chip,fill=white,draw=dGray!70,dashed,line width=0.5pt,text=dGray},
  fwd/.style={->,dTealLine,line width=0.8pt},
  tgt/.style={->,dOrange,line width=0.9pt},
  dual/.style={->,dCoral,dashed,line width=0.9pt},
  dep/.style={->,dTealLine,line width=0.8pt},
  pdep/.style={->,dGray,dashed,line width=0.7pt},
  nn/.style={circle,minimum size=0.26cm,inner sep=0pt,draw=dInk!70,line width=0.35pt},
}

% Lane frame in the paper style: two shadow copies behind a teal rounded frame.
\newcommand{\lane}[4]{% #1 x0, #2 y0, #3 x1, #4 y1
  \foreach \o/\c in {0.16/dShadow!60,0.08/dShadow}
    \draw[\c,fill=white,rounded corners=5pt,line width=0.6pt]
      (#1+\o,#2-\o) rectangle (#3+\o,#4-\o);
  \draw[dTealLine,fill=white,rounded corners=5pt,line width=0.9pt] (#1,#2) rectangle (#3,#4);}

% Plain card frame (no shadow copies), for problems and losses.
\newcommand{\cardbox}[4]{\draw[dShadow,fill=white,rounded corners=4pt,line width=0.6pt] (#1,#2) rectangle (#3,#4);}

% Network glyph on one sheet. #1 = centre, #2 = name prefix, #3 = number of stacked
% sheets behind (0, 1 or 2).
\newcommand{\net}[3]{%
  \begin{scope}[shift={(#1)}]
    \ifnum#3>1 \node[sheet,minimum width=1.9cm,minimum height=1.8cm] at (0.30,0.22) {};\fi
    \ifnum#3>0 \node[sheet,minimum width=1.9cm,minimum height=1.8cm] at (0.15,0.11) {};\fi
    \node[sheet,minimum width=1.9cm,minimum height=1.8cm,fill=white] (#2sheet) at (0,0) {};
    \foreach \y [count=\i] in {0.36,0,-0.36} \node[nn,fill=dGray!45] (#2in\i) at (-0.64,\y) {};
    \foreach \y [count=\i] in {0.58,0.29,0,-0.29,-0.58} \node[nn,fill=dNode] (#2ha\i) at (-0.21,\y) {};
    \foreach \y [count=\i] in {0.44,0.15,-0.15,-0.44} \node[nn,fill=dNode] (#2hb\i) at (0.21,\y) {};
    \node[nn,fill=dOrange] (#2out) at (0.64,0) {};
    \foreach \a in {1,...,3} \foreach \b in {1,...,5} \draw[dGray!60,line width=0.2pt] (#2in\a) -- (#2ha\b);
    \foreach \a in {1,...,5} \foreach \b in {1,...,4} \draw[dGray!60,line width=0.2pt] (#2ha\a) -- (#2hb\b);
    \foreach \a in {1,...,4} \draw[dGray!60,line width=0.2pt] (#2hb\a) -- (#2out);
  \end{scope}}

% Jacobian mask glyph: an n x d grid with a few entries selected. #1 = centre.
\newcommand{\maskgrid}[1]{%
  \begin{scope}[shift={(#1)}]
    \foreach \i in {0,...,4} \foreach \j in {0,...,3}
      \draw[dShadow,fill=white,line width=0.4pt] (\j*0.2-0.4,-\i*0.2+0.5) rectangle ++(0.2,-0.2);
    \foreach \i/\j in {0/1,1/3,2/0,3/2,4/1}
      \draw[dOrange,fill=dOrange,line width=0.4pt] (\j*0.2-0.4,-\i*0.2+0.5) rectangle ++(0.2,-0.2);
  \end{scope}}

% Code names inside italic notes stay upright.
\newcommand{\cmd}[1]{{\upshape\ttfamily #1}}

% Software-architecture helpers: package cards, type chips and chip rows.
\tikzset{
  pkg/.style={rectangle,rounded corners=4pt,fill=white,draw=dShadow,line width=0.6pt,
              align=center,inner xsep=5pt,inner ysep=4pt,font=\footnotesize},
  pkgcore/.style={pkg,draw=dTealLine,line width=0.9pt},
  pkgplan/.style={pkg,draw=dGray,dashed,text=dGray},
  tchip/.style={chip,font=\ttfamily\footnotesize,inner xsep=3pt,inner ysep=2.5pt},
  tplan/.style={tchip,fill=white,draw=dGray,dashed,line width=0.5pt,text=dGray},
  absname/.style={anchor=west,font=\ttfamily\footnotesize,text=dNavy,inner sep=1pt},
}
% \chiprow{x}{y}{style}{A,B,C}: chips left to right from (x,y), anchored west.
\newcommand{\chiprow}[4]{%
  \coordinate (cr@next) at (#1,#2);
  \foreach \t in {#4} {%
    \node[#3,anchor=west] (cr@last) at (cr@next) {\t};
    \coordinate (cr@next) at ($(cr@last.east)+(0.12,0)$);}}
